\subsection{Iterative Methods}

\begin{definition} \label{def:krylov-subspace} \hlt{Krylov Subspace}\\
For $A\in\operatorname{Mat}(n,n,\mathbb{F})$ and a vector $\mathbf{b}\in\mathbb{F}^n$, the $m$th \hlt{Krylov subspace} generated by $(A,\mathbf{b})$ is\\
$\mathcal{K}_m(A,\mathbf{b})\defeq\operatorname{span}\{\mathbf{b},A\mathbf{b},\ldots,A^{m-1}\mathbf{b}\}$.
\end{definition}

\begin{perspective} \label{per:krylov-methods-preview} \hlt{Krylov Methods Preview}\\
Krylov methods work in $\mathcal{K}_m(A,\mathbf{b})$: Arnoldi and Lanczos build bases, while CG and GMRES impose projection or residual-minimization conditions using only matrix-vector products.
\end{perspective}

\begin{method} \label{meth:arnoldi-lanczos-cg-preview} \hlt{Arnoldi, Lanczos, and CG Preview}\\
Let $\mathcal{K}_m=\mathcal{K}_m(A,\mathbf{b})$ and let $Q_m=[\mathbf{q}_1\ \cdots\ \mathbf{q}_m]$ be an orthonormal basis.
\begin{enumerate}[label=(\roman*)]
\setlength{\itemsep}{0pt}
\item \hlt{Arnoldi}: $AQ_m=Q_mH_m+h_{m+1,m}\mathbf{q}_{m+1}\mathbf{e}_m^t$, with $H_m$ upper Hessenberg.
\item \hlt{Lanczos}: if $A=A^*$, then $H_m$ is tridiagonal and $A\mathbf{q}_j=\beta_{j-1}\mathbf{q}_{j-1}+\alpha_j\mathbf{q}_j+\beta_j\mathbf{q}_{j+1}$.
\item \hlt{CG}: if $A\succ0$, choose $\mathbf{x}_k\in\mathbf{x}_0+\mathcal{K}_k(A,\mathbf{r}_0)$ with $\mathbf{r}_k\perp\mathcal{K}_k(A,\mathbf{r}_0)$; search directions satisfy $\mathbf{p}_i^*A\mathbf{p}_j=0$, $i\neq j$.
\end{enumerate}
Each step uses one matrix-vector product.
\end{method}

\begin{definition} \label{def:matrix-splitting-iteration} \hlt{Matrix Splitting Iteration}\\
Let $A=S-T$ with $S$ invertible, and set $B\defeq S^{-1}T$.\\
The \hlt{splitting iteration} for $A\mathbf{x}=\mathbf{b}$ is $S\mathbf{x}_{k+1}=T\mathbf{x}_k+\mathbf{b}$, and $\mathbf{x}_{k+1}=B\mathbf{x}_k+S^{-1}\mathbf{b}$.\\
For a solution $\mathbf{x}_*$, the error $\mathbf{e}_k\defeq\mathbf{x}_k-\mathbf{x}_*$ satisfies $\mathbf{e}_{k+1}=B\mathbf{e}_k$.
\end{definition}

\begin{theorem} \label{thm:splitting-iteration-convergence-spectral-radius} \hlt{Splitting Convergence Criterion}\\
Let $A=S-T$, with $S$ invertible, and set $B=S^{-1}T$. For $\mathbf{x}_{k+1}=B\mathbf{x}_k+S^{-1}\mathbf{b}$:
\begin{enumerate}[label=(\roman*)]
\setlength{\itemsep}{0pt}
\item if $A$ is invertible, then $\mathbf{x}_k\to A^{-1}\mathbf{b}$ for all $\mathbf{x}_0,\mathbf{b}$ if and only if $\rho(B)<1$;
\item $\rho(B)<1$ implies $A$ is invertible.
\end{enumerate}
\begin{proof}
\begin{enumerate}[label=(\roman*)]
\setlength{\itemsep}{0pt}
\item Let $\mathbf{x}_*=A^{-1}\mathbf{b}$ and $\mathbf{e}_k=\mathbf{x}_k-\mathbf{x}_*$. Then $\mathbf{e}_k=B^k\mathbf{e}_0$, so convergence for all $\mathbf{x}_0,\mathbf{b}$ is equivalent to $B^k\to0$, equivalently $\rho(B)<1$ by Theorem \ref{thm:matrix-powers-converge-to-zero-spectral-radius}.
\item If $\rho(B)<1$, then $1\notin\operatorname{spec}(B)$, so $I_n-B$ is invertible. Since $A=S-T=S(I_n-B)$, $A$ is invertible.
\end{enumerate}
\hlt{Structural Commentary}: A stationary method is a matrix-power problem in disguise.\\
\hlt{Key Insight}: The iteration matrix $B=S^{-1}T$ controls the errors.
\end{proof}
\end{theorem}

\begin{corollary} \label{cor:singular-matrix-splitting-obstruction} \hlt{Singular Matrices Obstruct Splitting Convergence}\\
Let $A=S-T$, with $S$ invertible. If $A$ is singular, then $1\in\operatorname{spec}(S^{-1}T)$.\\
Hence $\rho(S^{-1}T)\geq1$, so $\rho(S^{-1}T)<1$ cannot hold.
\begin{proof}
Choose $\mathbf{0}\neq\mathbf{x}\in\ker A$. Then $(S-T)\mathbf{x}=\mathbf{0}$, so $S\mathbf{x}=T\mathbf{x}$ and $S^{-1}T\mathbf{x}=\mathbf{x}$. Thus $1$ is eigenvalue of $S^{-1}T$.\\
\hlt{Structural Commentary}: Singular $A$ forces an undamped error mode.\\
\hlt{Key Insight}: A null vector of $A$ becomes a fixed error mode.
\end{proof}
\end{corollary}

\begin{method} \label{meth:stationary-iterative-method} \hlt{Stationary Iterative Method}\\
For $A\mathbf{x}=\mathbf{b}$, a stationary iteration is:
\begin{enumerate}[label=(\roman*)]
\setlength{\itemsep}{0pt}
\item Choose fixed $M$ and $\mathbf{c}$.
\item Update $\mathbf{x}_{k+1}=M\mathbf{x}_k+\mathbf{c}$.
\item If $\mathbf{x}_*$ exists, use $\mathbf{x}_{k+1}-\mathbf{x}_*=M(\mathbf{x}_k-\mathbf{x}_*)$ and Theorem \ref{thm:matrix-powers-converge-to-zero-spectral-radius} for convergence.
\end{enumerate}
\end{method}

\begin{method} \label{meth:jacobi-gauss-seidel-sor-splittings} \hlt{Jacobi, Gauss--Seidel, and SOR Splittings}\\
Let $A=D+L+U$, with $D$ diagonal, $L$ strictly lower triangular, and $U$ strictly upper triangular. For a splitting $A=S-T$, use $S\mathbf{x}_{k+1}=T\mathbf{x}_k+\mathbf{b}$ and $B=S^{-1}T$.
\begin{enumerate}[label=(\roman*)]
\setlength{\itemsep}{0pt}
\item \hlt{Jacobi}: $S=D$, $T=-(L+U)$, $B_J=-D^{-1}(L+U)$.
\item \hlt{Gauss--Seidel}: $S=D+L$, $T=-U$, $B_{\operatorname{GS}}=-(D+L)^{-1}U$.
\item \hlt{SOR}: for $\omega\neq0$, $S_\omega=\frac1\omega D+L$; $T_\omega=\left(\frac1\omega-1\right)D-U$; and $B_\omega=S_\omega^{-1}T_\omega$.
\end{enumerate}
Convergence is controlled by $\rho(B)$; $\omega=1$ gives Gauss--Seidel.
\end{method}

\begin{corollary} \label{cor:strict-diagonal-dominance-jacobi-convergence} \hlt{Strict Diagonal Dominance Gives Jacobi Convergence}\\
If $A=[a_{ij}]\in\operatorname{Mat}(n,n,\mathbb{C})$ is strictly row diagonally dominant, then Jacobi converges for every right-hand side.
\begin{proof}
Here $a_{ii}\neq0$ and $B_J=-D^{-1}(L+U)=I_n-D^{-1}A$. Thus $(B_J)_{ii}=0$ and $\sum\limits_{j\neq i}\abs{(B_J)_{ij}}=\sum\limits_{j\neq i}\abs{a_{ij}}/\abs{a_{ii}}<1$.\\
Gersgorin gives $\operatorname{spec}(B_J)\subset\{z:\abs{z}<1\}$, so $\rho(B_J)<1$; Theorem \ref{thm:splitting-iteration-convergence-spectral-radius} applies.\\
\hlt{Structural Commentary}: Diagonal dominance puts the Jacobi spectrum in the unit disk.\\
\hlt{Key Insight}: Small off-diagonal row sums force the error matrix to have spectral radius below $1$.
\end{proof}
\end{corollary}

\begin{example} \label{ex:two-by-two-jacobi-gauss-seidel-convergence} \hlt{Two-by-Two Jacobi and Gauss--Seidel}\\
For $A=\begin{bmatrix}2&1\\1&2\end{bmatrix}$, $\mathbf{b}\in\mathbb{R}^2$, write $A=D+L+U$.\\
The Jacobi iteration matrix is $B_J=-D^{-1}(L+U)= \begin{bmatrix} 0&-\frac12\\ -\frac12&0 \end{bmatrix}$, so $\rho(B_J)=\frac12$.\\
The Gauss--Seidel iteration matrix is $B_{\operatorname{GS}}=-(D+L)^{-1}U= \begin{bmatrix} 0&-\frac12\\ 0&\frac14 \end{bmatrix}$, so $\rho(B_{\operatorname{GS}})=\frac14$.\\
Both converge by Theorem \ref{thm:splitting-iteration-convergence-spectral-radius}, and smaller spectral radius predicts faster asymptotic decay for Gauss--Seidel.
\end{example}

\begin{example} \label{ex:poisson-tridiagonal-stationary-rates} \hlt{Poisson Tridiagonal Stationary Rates}\\
For $A_n=\operatorname{tridiag}(-1,2,-1)$, Jacobi gives $B_J=I_n-\frac12A_n$; $\lambda_k(B_J)=\cos\frac{k\pi}{n+1}$; and $\rho(B_J)=\cos\frac{\pi}{n+1}$.\\
For the same ordering, Gauss--Seidel satisfies $\rho(B_{\operatorname{GS}})=\rho(B_J)^2=\cos^2\frac{\pi}{n+1}$.
As $n\to\infty$, both radii tend to $1$; fine grids need acceleration or preconditioning.
\end{example}

\begin{perspective} \label{per:preconditioning-preview} \hlt{Preconditioning Preview}\\
A preconditioner replaces $A\mathbf{x}=\mathbf{b}$ by an equivalent system such as $M^{-1}A\mathbf{x}=M^{-1}\mathbf{b}$ or $AM^{-1}\mathbf{y}=\mathbf{b}$, with $\mathbf{x}=M^{-1}\mathbf{y}$. It should improve conditioning or convergence while being much cheaper to apply than $A^{-1}$.
\end{perspective}
